\begin{table}[htbp]
\centering
\footnotesize
\caption{Problema D, meia-lua, Passos 2 a 4 (openai/gpt-4o-mini): zero-shot versus few-shot por espaço.}
\label{tab:bloco1_D_passo2_vs_passo3_gpt4mini}
\begin{tabular}{llrrrrrrrrr}
\toprule
 &  & \multicolumn{3}{c}{Fidelidade (vs LLM)} & \multicolumn{3}{c}{Acur.\ da métrica vs real} & \multicolumn{3}{c}{Acur.\ do LLM vs real} \\ \cmidrule(lr){3-5} \cmidrule(lr){6-8} \cmidrule(lr){9-11}
Espaço & Few-shot & zero & few & $\Delta$ & zero & few & $\Delta$ & zero & few & $\Delta$ \\
\midrule
R$^2$ & 4-shot & — & 89,8 (76,7–97,3) & — & — & 82,9 (76,7–86,7) & — & — & 81,9 (68,5–92,5) & — \\
R$^2$ & 10-shot & — & 92,3 (81,4–97,1) & — & — & 85,1 (77,3–88,0) & — & — & 89,2 (83,6–96,4) & — \\
R$^2$ & 20-shot & — & 93,8 (92,3–96,2) & — & — & 86,1 (83,3–88,0) & — & — & 90,4 (86,9–95,4) & — \\
R$^2$ & 40-shot & — & 92,1 (88,2–94,5) & — & — & 86,1 (84,0–87,3) & — & — & 91,6 (86,4–94,5) & — \\
\midrule
R$^3$ & 4-shot & — & 90,3 (78,8–97,3) & — & — & 82,2 (78,0–86,7) & — & — & 81,9 (68,5–92,5) & — \\
R$^3$ & 10-shot & — & 92,5 (81,4–97,1) & — & — & 85,3 (78,0–87,3) & — & — & 89,2 (83,6–96,4) & — \\
R$^3$ & 20-shot & — & 93,5 (91,5–96,2) & — & — & 85,6 (82,7–87,3) & — & — & 90,4 (86,9–95,4) & — \\
R$^3$ & 40-shot & — & 92,1 (90,0–94,5) & — & — & 85,6 (84,0–86,7) & — & — & 91,6 (86,4–94,5) & — \\
\midrule
R$^4$ & 4-shot & — & 94,2 (89,7–99,3) & — & — & 81,7 (75,3–86,7) & — & — & 81,9 (68,5–92,5) & — \\
R$^4$ & 10-shot & — & 93,6 (80,0–98,6) & — & — & 85,7 (80,7–88,0) & — & — & 89,2 (83,6–96,4) & — \\
R$^4$ & 20-shot & — & 94,7 (93,1–97,7) & — & — & 86,5 (84,7–88,0) & — & — & 90,4 (86,9–95,4) & — \\
R$^4$ & 40-shot & — & 93,1 (90,0–95,5) & — & — & 86,7 (84,7–88,7) & — & — & 91,6 (86,4–94,5) & — \\
\bottomrule
\multicolumn{11}{p{0.95\linewidth}}{\footnotesize Meia-lua, Perceptron. Fidelidade (vs LLM); acurácia (vs rótulo real) da métrica e do LLM como contexto. Colapsos (e quase-colapsos) do zero-shot não entram em nenhuma média, nem na acurácia do LLM (``—'' quando todas as sementes colapsaram; detalhe na bloco1\_table\_10). Em zero-shot (e com nomes de classe neutros) a correspondência entre nome e classe é arbitrária: a acurácia (vs rótulo real) do LLM pode sair abaixo de 50\%; o valor invariante à permutação é max($a$, 1 $-$ $a$). $\Delta$ = few-shot $-$ zero-shot. Passos do Problema D: zero-shot = Passo 2 (o LLM rotula sem exemplos); $k$-shot = Passo 3 (o LLM rotula com $k$ exemplos do rótulo real); a métrica de cada linha é estimada no Passo 4; teto = Passo 5 (W ajustado diretamente ao rótulo real).} \\
\end{tabular}
\end{table}
